\documentclass[9pt]{beamer}
\mode<presentation>

% Theme choice:
\usetheme{Madrid}%Darmstadt 
\usepackage{xcolor}
\definecolor{mGreen}{rgb}{0,0.6,0}
\definecolor{mRed}{rgb}{0.6,0,0}
\usecolortheme[RGB={0,100,50}]{structure}
 \usefonttheme{structurebold} 
\setbeamercovered{invisible}
%\setbeamertemplate{navigation symbols}{}
\usepackage{dynblocks}
\usepackage{textpos}
\usepackage{amsfonts}
\usepackage{amsmath}
\usepackage{amssymb}
\usepackage{amsthm}
\usepackage{listings}
\usepackage{tikz}
\usepackage{pgfplots}
\pgfplotsset{compat=1.18}
\newtheorem{teorema}{Teorema}\usepackage{mathtools}
\renewcommand{\proofname}{Prova}
\usepackage[portuguese]{babel}

\uselanguage{Portuguese}
\languagepath{Portuguese}

\deftranslation[to=Portuguese]{Theorem}{Teorema}
\deftranslation[to=Portuguese]{Lemma}{Lema}
\deftranslation[to=Portuguese]{Corollary}{Corolário}
\deftranslation[to=Portuguese]{Proof}{Prova}

\usepackage{dutchcal}
\usepackage{hyperref}
\usepackage[utf8]{inputenc}
\usepackage[T1]{fontenc}
\usepackage{amsthm}
\usepackage{tikz}
\usetikzlibrary{shapes,arrows,positioning}

\usepackage[affil-it]{authblk} 
\usepackage{etoolbox}
\usepackage{lmodern}

\makeatletter
\patchcmd{\@maketitle}{\LARGE \@title}{\fontsize{18}{19.2}\selectfont\@title}{}{}
\makeatother

\renewcommand\Authfont{\fontsize{12}{14.4}\selectfont}
\renewcommand\Affilfont{\fontsize{9}{10.8}\itshape}

% Title page details: 
\title[Presentation]{ENGENHARIA DA QUALIDADE E CONFIABILIDADE} %add title


\institute[UniRV]{\textbf{\large  \textcolor{structure}{\textbf{\large
Métricas Google SRE
}} %add name 
% \hspace{0.5pt}
\\ \qquad
\\Professor: Me. Sergio Souza Novak 
\\Faculdade de Engenharia de Software}
} 
% \titlegraphic{\includegraphics[width=5cm]{msu.eps}}
\titlegraphic{\includegraphics[width=5cm]{Universidade_de_Rio_Verde.jpg}}
\date{}

% Logo only on title page
%\logo{\includegraphics[width=1cm,keepaspectratio]{Logo.png}}
\usepackage[pscoord]{eso-pic}


\begin{document}

% Title page frame
%\begin{frame}
    %\titlepage 
%\end{frame}
\frame[plain]{\titlepage}
%logo
\AddToShipoutPictureFG{
    \put(\LenToUnit{.92\paperwidth},
    \LenToUnit{.185\paperheight})
    {\vtop{{\null}
    \makebox{\includegraphics[width=.8cm,keepaspectratio]{Universidade_de_Rio_Verde.jpg}}}}}


\begin{frame}{Sumário}
  \tableofcontents
\end{frame}

% Slide 1
\section{Introdução ás métricas}
\begin{frame}{Mensurando os Riscos a um Serviço}

\begin{block}{Prática-Padrão}
Como prática-padrão no Google, busca-se definir uma \textbf{métrica objetiva}
para representar a propriedade que queremos otimizar.
\end{block}

\begin{block}{Por que usar métricas?}
\begin{itemize}
    \item Avaliar o desempenho atual
    \item Monitorar melhorias
    \item Identificar degradações ao longo do tempo
\end{itemize}
\end{block}

\begin{alertblock}{Desafio}
Nem sempre é simples reduzir todos os riscos a uma única métrica.
\end{alertblock}

\end{frame}
% Slide 2
\begin{frame}{Impactos das Falhas de Serviço}

\begin{block}{Consequências Potenciais}
\begin{itemize}
    \item Insatisfação dos usuários
    \item Perda de confiança
    \item Perda de receita
    \item Impacto na marca
    \item Cobertura negativa da imprensa
\end{itemize}
\end{block}

\begin{center}
\begin{tikzpicture}[
    node distance=1.8cm,
    every node/.style={draw, rounded corners, align=center}
]
\node (falha) {Falha de\\ Serviço};
\node (usuarios) [below left=of falha] {Usuários};
\node (receita) [below=of falha] {Receita};
\node (marca) [below right=of falha] {Marca};

\draw[->] (falha) -- (usuarios);
\draw[->] (falha) -- (receita);
\draw[->] (falha) -- (marca);
\end{tikzpicture}
\end{center}

\end{frame}
% Slide 3
\section{Downtime Não Planejado}

\begin{frame}{Foco no Downtime Não Planejado}

\begin{block}{Por que focar no Downtime?}
Ao lidar com riscos em sistemas complexos, existem muitos impactos possíveis 
(financeiros, reputacionais, operacionais). Para garantir \textbf{padronização} 
e comparabilidade entre diferentes serviços, escolhe-se uma métrica comum:
\textbf{o downtime não planejado}.
\end{block}

\begin{exampleblock}{Definição Formal}
Downtime não planejado corresponde aos períodos em que o serviço
fica indisponível sem que essa interrupção tenha sido previamente
programada ou comunicada.
\end{exampleblock}

\begin{alertblock}{Importância}
\begin{itemize}
    \item Afeta diretamente a experiência do usuário
    \item Impacta SLAs e contratos
    \item Pode gerar perdas financeiras
    \item Compromete a confiança no sistema
\end{itemize}
\end{alertblock}

\begin{center}
\begin{tikzpicture}
\draw[->] (0,0) -- (8,0) node[right]{Tempo};

\draw[thick] (0,0.5) -- (2,0.5);
\draw[thick, red] (2,0.5) -- (3,0.5);
\draw[thick] (3,0.5) -- (6,0.5);
\draw[thick, red] (6,0.5) -- (7,0.5);
\draw[thick] (7,0.5) -- (8,0.5);

\node at (1,0.8) {\small Uptime};
\node at (2.5,1) {\small Downtime};
\node at (4.5,0.8) {\small Uptime};
\node at (6.5,1) {\small Downtime};
\node at (7.5,0.8) {\small Uptime};
\end{tikzpicture}
\end{center}

\end{frame}
\section{Disponibilidade do Serviço}

\begin{frame}{Disponibilidade do Serviço}

\begin{block}{Conceito de Disponibilidade}
Disponibilidade representa a probabilidade de um sistema estar
operacional e acessível quando requisitado pelo usuário.
Em termos práticos, expressa o nível tolerável de downtime não planejado.
\end{block}

\begin{block}{Expressão em "Noves"}
A disponibilidade é tradicionalmente representada pelo número de "noves":

\begin{itemize}
    \item 99,9\%  → até 8h45min de indisponibilidade por ano
    \item 99,99\% → até 52min por ano
    \item 99,999\% → cerca de 5min por ano
\end{itemize}

Cada nove adicional reduz o downtime em aproximadamente
\textbf{uma ordem de magnitude}.
\end{block}

\begin{alertblock}{Interpretação}
Pequenas diferenças percentuais representam grandes diferenças
no tempo total de indisponibilidade anual.
\end{alertblock}

\begin{center}
\begin{tikzpicture}
\draw[->] (0,0) -- (6,0) node[right]{Disponibilidade};

\draw (1,0.1) -- (1,-0.1);
\draw (3,0.1) -- (3,-0.1);
\draw (5,0.1) -- (5,-0.1);

\node[below] at (1,-0.1) {99,9\%};
\node[below] at (3,-0.1) {99,99\%};
\node[below] at (5,-0.1) {99,999\%};

\node[above] at (1,0.2) {\tiny ~8h/ano};
\node[above] at (3,0.2) {\tiny ~52min/ano};
\node[above] at (5,0.2) {\tiny ~5min/ano};
\end{tikzpicture}
\end{center}

\end{frame} 
% Slide 5
\begin{frame}{Cálculo Tradicional da Disponibilidade}

\begin{block}{Equação 3.1}
\[
\text{Disponibilidade} =
\frac{\text{Tempo de Uptime}}{\text{Tempo Total}}
\]
\end{block}

\begin{block}{Interpretação}
\begin{itemize}
    \item Mede proporção de tempo operacional
    \item Base da engenharia de confiabilidade
\end{itemize}
\end{block}

\begin{center}
\begin{tikzpicture}[node distance=2cm]
\node (total) [draw, rounded corners] {Tempo Total};
\node (uptime) [draw, rounded corners, below left=of total] {Uptime};
\node (downtime) [draw, rounded corners, below right=of total] {Downtime};

\draw[->] (total) -- (uptime);
\draw[->] (total) -- (downtime);
\end{tikzpicture}
\end{center}

\end{frame}
% Slide 6
\begin{frame}{Resumo}

\begin{block}{Principais Pontos}
\begin{itemize}
    \item Riscos devem ser traduzidos em métricas
    \item Foco no downtime não planejado
    \item Disponibilidade expressa em “noves”
    \item Métricas permitem melhoria contínua
\end{itemize}
\end{block}

\begin{center}
\begin{tikzpicture}[node distance=1.5cm]
\node (riscos) [draw, rounded corners] {Riscos};
\node (metricas) [draw, rounded corners, below=of riscos] {Métricas};
\node (melhoria) [draw, rounded corners, below=of metricas] {Melhoria Contínua};

\draw[->] (riscos) -- (metricas);
\draw[->] (metricas) -- (melhoria);
\end{tikzpicture}
\end{center}

\end{frame}
\begin{frame}{Por que Métricas Geram Melhoria Contínua?}

\begin{block}{1. Tornam o desempenho observável}
O que não é medido permanece no campo da percepção.
Métricas transformam impressões subjetivas em dados objetivos,
permitindo análise técnica e comparação temporal.
\end{block}

\begin{block}{2. Criam ciclos de feedback}
A melhoria contínua depende de realimentação.
Ao medir desempenho antes e depois de uma intervenção,
é possível avaliar impacto, validar hipóteses e corrigir estratégias.
\end{block}

\begin{block}{3. Permitem definição de metas claras}
Indicadores quantitativos viabilizam:
\begin{itemize}
    \item Estabelecimento de SLAs
    \item Comparação com benchmarks
    \item Priorização de investimentos
\end{itemize}
Sem métricas, decisões tendem a ser reativas.
\end{block}

\end{frame}
\begin{frame}{O Mecanismo da Melhoria Contínua Baseada em Métricas}

\begin{block}{Ciclo Estruturado de Aprendizado}
A medição sistemática viabiliza ciclos iterativos como o PDCA:

\begin{itemize}
    \item \textbf{Plan} – definir metas (ex: reduzir MTTR)
    \item \textbf{Do} – implementar melhorias
    \item \textbf{Check} – medir resultados
    \item \textbf{Act} – ajustar processos
\end{itemize}

Sem a etapa de medição, o ciclo se interrompe.
\end{block}

\begin{block}{Efeitos Organizacionais}
Métricas adequadas:
\begin{itemize}
    \item Alinham incentivos e prioridades
    \item Reduzem ambiguidade na tomada de decisão
    \item Permitem identificar tendências (melhoria ou degradação)
    \item Transformam incidentes em aprendizado cumulativo
\end{itemize}
\end{block}
\end{frame}
\section{Disponibilidade em sistemas distribuídos}

\begin{frame}{Problema da Disponibilidade Baseada em Tempo em Sistemas Distribuídos}

\begin{block}{Contexto em Sistemas Distribuídos Globais}

\begin{itemize}
    \item Múltiplos data centers geograficamente distribuídos
    \item Isolamento regional de falhas
    \item Impacto parcial de tráfego
    \item Sistema quase nunca 100\% down
\end{itemize}

\end{block}

\begin{alertblock}{Implicação}
Medir disponibilidade apenas como “tempo totalmente fora do ar”
torna-se pouco representativo, pois o sistema pode estar parcialmente
operacional durante toda a falha.
\end{alertblock}

\begin{center}
\begin{tikzpicture}[
    node distance=1.8cm,
    every node/.style={font=\small}
]

% Usuários
\node (u1) at (0,2) {Usuários - América};
\node (u2) at (0,1) {Usuários - Europa};
\node (u3) at (0,0) {Usuários - Ásia};

% Data centers
\node[draw, rounded corners, minimum width=2.5cm] (dc1) at (6,2) {Data Center A};
\node[draw, rounded corners, minimum width=2.5cm, fill=red!20] (dc2) at (6,1) {Data Center B (Falha)};
\node[draw, rounded corners, minimum width=2.5cm] (dc3) at (6,0) {Data Center C};

% Conexões funcionando
\draw[->] (u1) -- (dc1);
\draw[->] (u3) -- (dc3);

% Conexão com falha
\draw[->, dashed, red] (u2) -- (dc2);

% Observação
\node at (3,-0.8) {\footnotesize Sistema parcialmente operacional};

\end{tikzpicture}
\end{center}

\end{frame} 
\section{O que é uma API REST?}

\begin{frame}{O que é uma API REST?}

\begin{block}{Definição}
API (Application Programming Interface) é um conjunto de regras
que permite que sistemas se comuniquem entre si.

REST (Representational State Transfer) é um estilo arquitetural
baseado em recursos acessados via HTTP.
\end{block}

\begin{block}{Características Principais}
\begin{itemize}
    \item Comunicação cliente-servidor
    \item Uso de métodos HTTP (GET, POST, PUT, DELETE)
    \item Recursos identificados por URLs
    \item Respostas normalmente em JSON
\end{itemize}
\end{block}

\begin{center}
\begin{tikzpicture}

\node[draw, rounded corners, minimum width=2.5cm, minimum height=1cm]
(cliente) at (0,0) {Cliente};

\node[draw, rounded corners, minimum width=3cm, minimum height=1cm]
(api) at (8,0) {API REST};

\draw[->] (cliente) -- node[above]{HTTP Request} (api);
\draw[<-] (cliente) -- node[below]{HTTP Response (JSON)} (api);

\end{tikzpicture}
\end{center}

\end{frame}
\section{O que é uma requisição?}

\begin{frame}{O que é uma Requisição?}

\begin{block}{Definição}
Uma requisição é uma mensagem enviada por um cliente
a um servidor solicitando a execução de uma operação
sobre um recurso.
\end{block}

\begin{block}{Componentes de uma Requisição HTTP}
\begin{itemize}
    \item Método (GET, POST, PUT, DELETE)
    \item URL (recurso desejado)
    \item Cabeçalhos (headers)
    \item Corpo (body) — opcional
\end{itemize}
\end{block}

\begin{center}
\begin{tikzpicture}

\node[draw, rounded corners, align=left, minimum width=7cm]
(req) at (0,0) {
\texttt{POST /usuarios} \\
Host: exemplo.com \\
Content-Type: application/json \\
\{ "nome": "Ana" \}
};

\end{tikzpicture}
\end{center}

\end{frame}
% =========================
% SLIDE GET
% =========================
\section{Métodos HTTP}
\subsection{GET}

\begin{frame}{Método HTTP: GET}

\begin{block}{O que é?}
Obtém dados de um recurso no servidor.
Não altera estado. Seguro e idempotente.
\end{block}

\begin{block}{Exemplo de Requisição}
\ttfamily\footnotesize
GET /usuarios/10 HTTP/1.1 \\
Host: api.exemplo.com \\
Accept: application/json
\end{block}

\begin{block}{Exemplo de Resposta}
\ttfamily\footnotesize
HTTP/1.1 200 OK \\
Content-Type: application/json \\

\{ \\
\quad "id": 10, \\
\quad "nome": "Ana", \\
\quad "email": "ana@email.com" \\
\}
\end{block}

\begin{center}
\begin{tikzpicture}
\node[draw, rounded corners] (c) at (0,0) {Cliente};
\node[draw, rounded corners] (s) at (5,0) {Servidor};
\draw[->] (c) -- node[above]{GET} (s);
\draw[<-] (c) -- node[below]{200 OK} (s);
\end{tikzpicture}
\end{center}

\end{frame}

% =========================
% SLIDE POST
% =========================
\subsection{POST}

\begin{frame}{Método HTTP: POST}

\begin{block}{O que é?}
Cria um novo recurso no servidor.
Não é idempotente.
\end{block}

\begin{block}{Exemplo de Requisição}
\ttfamily\footnotesize
POST /usuarios HTTP/1.1 \\
Host: api.exemplo.com \\
Content-Type: application/json \\

\{ \\
\quad "nome": "Carlos", \\
\quad "email": "carlos@email.com" \\
\}
\end{block}

\begin{block}{Exemplo de Resposta}
\ttfamily\footnotesize
HTTP/1.1 201 Created \\
Location: /usuarios/11
\end{block}

\begin{center}
\begin{tikzpicture}
\node[draw, rounded corners] (c) at (0,0) {Cliente};
\node[draw, rounded corners] (s) at (5,0) {Servidor};
\draw[->] (c) -- node[above]{POST} (s);
\draw[<-] (c) -- node[below]{201 Created} (s);
\end{tikzpicture}
\end{center}

\end{frame}

% =========================
% SLIDE PUT
% =========================
\subsection{PUT}

\begin{frame}{Método HTTP: PUT}

\begin{block}{O que é?}
Atualiza completamente um recurso existente.
É idempotente.
\end{block}

\begin{block}{Exemplo de Requisição}
\ttfamily\footnotesize
PUT /usuarios/10 HTTP/1.1 \\
Host: api.exemplo.com \\
Content-Type: application/json \\

\{ \\
\quad "nome": "Ana Souza", \\
\quad "email": "ana@email.com" \\
\}
\end{block}

\begin{block}{Exemplo de Resposta}
\ttfamily\footnotesize
HTTP/1.1 200 OK
\end{block}

\begin{center}
\begin{tikzpicture}
\node[draw, rounded corners] (c) at (0,0) {Cliente};
\node[draw, rounded corners] (s) at (5,0) {Servidor};
\draw[->] (c) -- node[above]{PUT} (s);
\draw[<-] (c) -- node[below]{200 OK} (s);
\end{tikzpicture}
\end{center}

\end{frame}

% =========================
% SLIDE DELETE
% =========================
\subsection{DELETE}

\begin{frame}{Método HTTP: DELETE}

\begin{block}{O que é?}
Remove um recurso do servidor.
É idempotente.
\end{block}

\begin{block}{Exemplo de Requisição}
\ttfamily\footnotesize
DELETE /usuarios/10 HTTP/1.1 \\
Host: api.exemplo.com
\end{block}

\begin{block}{Exemplo de Resposta}
\ttfamily\footnotesize
HTTP/1.1 204 No Content
\end{block}

\begin{center}
\begin{tikzpicture}
\node[draw, rounded corners] (c) at (0,0) {Cliente};
\node[draw, rounded corners] (s) at (5,0) {Servidor};
\draw[->] (c) -- node[above]{DELETE} (s);
\draw[<-] (c) -- node[below]{204 No Content} (s);
\end{tikzpicture}
\end{center}

\end{frame}

\section{Disponibilidade Baseada em Requisições}

\begin{frame}{Disponibilidade Baseada em Requisições}

\begin{block}{Definição Utilizada pelo Google}
Em sistemas distribuídos globais, a disponibilidade é definida como:

\[
Disponibilidade =
\frac{\text{Requisições bem sucedidas}}
{\text{Total de requisições}}
\]

A métrica é calculada em uma janela móvel
(ex: 1 dia, 1 semana ou 1 trimestre).
\end{block}

\begin{center}
\begin{tikzpicture}

% Total box
\node[draw, rounded corners, minimum width=6cm, minimum height=1cm]
(total) at (0,1.5) {Total de Requisições};

% Success portion
\node[draw, rounded corners, fill=green!20,
minimum width=4.5cm, minimum height=1cm]
(success) at (-0.75,0) {Requisições Bem Sucedidas};

% Failed portion
\node[draw, rounded corners, fill=red!20,
minimum width=1.2cm, minimum height=1cm]
(fail) at (2.7,0) {Erros};

\draw[->] (total) -- (success);
\draw[->] (total) -- (fail);

\end{tikzpicture}
\end{center}

\end{frame}
\begin{frame}{Exemplo de Cálculo Diário}

\begin{block}{Cenário}
\begin{itemize}
    \item 2.500.000 requisições em um dia
    \item Meta de disponibilidade: 99,99\%
\end{itemize}
\end{block}

\begin{block}{Erro Máximo Permitido}
\[
0,01\% \times 2.500.000 = 250 \text{ erros}
\]

O sistema pode falhar até 250 requisições
e ainda atingir a meta diária.
\end{block}

\begin{center}
\begin{tikzpicture}

\draw[fill=green!20] (0,0) rectangle (6,0.6);
\draw[fill=red!40] (6,0) rectangle (6.06,0.6);

\node at (3,0.3) {\small 2.499.750 Sucessos};
\node[above] at (6.03,0.6) {\tiny 250 Erros};

\end{tikzpicture}
\end{center}

\end{frame}
\begin{frame}{Por Que Medir por Requisições?}

\begin{block}{Perspectiva do Usuário}
Disponibilidade, na prática, significa:

\begin{center}
\textit{"Minha requisição funcionou ou não?"}
\end{center}

O usuário não percebe horas de uptime.
Ele percebe sucesso ou falha da ação executada.
\end{block}

\begin{block}{Vantagens do Modelo}
\begin{itemize}
    \item Reflete experiência real do usuário
    \item Funciona em sistemas parcialmente operacionais
    \item É diretamente mensurável por engenharia
    \item Permite metas quantitativas claras
\end{itemize}
\end{block}

\begin{center}
\begin{tikzpicture}

\node[draw, rounded corners] (user) at (0,1) {Usuário};
\node[draw, rounded corners] (req) at (3,1) {Requisição};
\node[draw, rounded corners] (resp) at (6,1) {Resposta};

\draw[->] (user) -- (req);
\draw[->] (req) -- (resp);

\node at (6,-0.2) {\small Sucesso ou Falha};

\end{tikzpicture}
\end{center}

\end{frame}
\begin{frame}[fragile]{Status 100 – Continue (Faixa 1xx – Informacional)}

\begin{block}{O que significa a faixa 1xx?}
A faixa 1xx representa respostas informativas.
O servidor recebeu a requisição e está processando,
mas ainda não concluiu a resposta final.
\end{block}

\begin{block}{Sobre o 100 Continue}
O status 100 indica que o cliente pode continuar enviando
o corpo da requisição. É comum em requisições POST grandes.
\end{block}

\textbf{Exemplo comum: envio de relatório grande}

\begin{block}{Request}
\begin{verbatim}
POST /caldeira/relatorio HTTP/1.1
Host: api.siderurgia.com
Expect: 100-continue
Content-Length: 500000
\end{verbatim}
\end{block}

\begin{block}{Resposta intermediária}
\begin{verbatim}
HTTP/1.1 100 Continue
\end{verbatim}
\end{block}

\begin{block}{Interpretação do QA}
Não é erro.
É apenas uma confirmação de que o servidor
está pronto para receber o corpo da requisição.
\end{block}

\end{frame}

\begin{frame}[fragile]{Status 200 – OK (Faixa 2xx – Sucesso)}

\begin{block}{O que significa a faixa 2xx?}
Indica que a requisição foi recebida,
processada e concluída com sucesso.
\end{block}

\begin{block}{Sobre o 200 OK}
É o código mais comum.
Significa que a operação foi executada corretamente.
\end{block}

\textbf{Exemplo frequente: consulta de status}

\begin{block}{Request}
\begin{verbatim}
GET /caldeira/status HTTP/1.1
Host: api.siderurgia.com
\end{verbatim}
\end{block}

\begin{block}{Response}
\begin{verbatim}
HTTP/1.1 200 OK
Content-Type: application/json

{
  "temperatura": 850,
  "pressao": 120,
  "status": "OPERANDO_NORMAL"
}
\end{verbatim}
\end{block}

\begin{block}{Interpretação do QA}
Fluxo feliz.
A operação foi concluída corretamente
e o sistema respondeu conforme esperado.
\end{block}

\end{frame}
\begin{frame}[fragile]{Status 300 – Redirecionamento (Faixa 3xx)}

\begin{block}{O que significa a faixa 3xx?}
Indica que o recurso solicitado
está disponível em outro local
ou exige ação adicional do cliente.
\end{block}

\begin{block}{Exemplo: 301 Moved Permanently}
O endpoint foi alterado permanentemente.
\end{block}

\begin{block}{Request}
\begin{verbatim}
GET /caldeira/dados HTTP/1.1
Host: api.siderurgia.com
\end{verbatim}
\end{block}

\begin{block}{Response}
\begin{verbatim}
HTTP/1.1 301 Moved Permanently
Location: /caldeira/status
\end{verbatim}
\end{block}

\begin{block}{Interpretação do QA}
Não é erro.
O cliente deve redirecionar automaticamente
para o novo endpoint indicado.
\end{block}

\end{frame}

\begin{frame}[fragile]{Status 404 – Not Found (Faixa 4xx – Erro do Cliente)}

\begin{block}{O que significa a faixa 4xx?}
Erros causados por requisições inválidas,
malformadas ou com recurso inexistente.
A responsabilidade é do cliente.
\end{block}

\begin{block}{Sobre o 404}
Indica que o recurso solicitado não existe.
\end{block}

\begin{block}{Request}
\begin{verbatim}
GET /caldeira/manutencao/9999 HTTP/1.1
Host: api.siderurgia.com
\end{verbatim}
\end{block}

\begin{block}{Response}
\begin{verbatim}
HTTP/1.1 404 Not Found

{
  "erro": "Registro não encontrado"
}
\end{verbatim}
\end{block}


\end{frame}

\begin{frame}[fragile]{Status 404 – Not Found (Faixa 4xx – Erro do Cliente)}


\begin{block}{Request}
\begin{verbatim}
GET /caldeira/manutencao/9999 HTTP/1.1
Host: api.siderurgia.com
\end{verbatim}
\end{block}

\begin{block}{Response}
\begin{verbatim}
HTTP/1.1 404 Not Found

{
  "erro": "Registro não encontrado"
}
\end{verbatim}
\end{block}

\begin{block}{Interpretação do QA}
Erro esperado.
Sistema deve responder corretamente
sem gerar falha interna.
\end{block}

\end{frame}
\begin{frame}[fragile]{Status 500 – Internal Server Error (Faixa 5xx)}
\begin{block}{O que significa a faixa 5xx?}
Indica falha interna no servidor.
O problema não está na requisição,
mas no processamento.
\end{block}
\begin{block}{Request}
\begin{verbatim}
POST /caldeira/manutencao HTTP/1.1
Content-Type: application/json

{
  "tecnico": "Carlos Silva"
}
\end{verbatim}
\end{block}

\begin{block}{Response}
\begin{verbatim}
HTTP/1.1 500 Internal Server Error

{
  "erro": "Falha interna no sistema"
}
\end{verbatim}
\end{block}

\begin{block}{Interpretação do QA}
Erro crítico.
Impacta disponibilidade.
Deve gerar incidente ou alerta.
\end{block}

\end{frame}

\begin{frame}[fragile]{Status 500 – Internal Server Error (Faixa 5xx)}

\begin{block}{Request}
\begin{verbatim}
POST /caldeira/manutencao HTTP/1.1
Content-Type: application/json

{
  "tecnico": "Carlos Silva"
}
\end{verbatim}
\end{block}

\begin{block}{Response}
\begin{verbatim}
HTTP/1.1 500 Internal Server Error

{
  "erro": "Falha interna no sistema"
}
\end{verbatim}
\end{block}

\begin{block}{Interpretação do QA}
Erro crítico.
Impacta disponibilidade.
Deve gerar incidente ou alerta.
\end{block}

\end{frame}
\begin{frame}
\centering
\vspace{1cm}

{\Huge \textbf{Estudo de Caso}} \\[0.5cm]
{\Large Teste de uso de sistema}

\vspace{1cm}

Sistema de Monitoramento de Caldeira Industrial \\
Siderurgia

\vfill
\small Exercício de Teste Manual de QA
\end{frame}


\begin{frame}[fragile]{GET – Consulta de Status da Caldeira}

\begin{block}{O que o QA fez}
O QA acessou o endpoint de status para verificar se a caldeira
está operando dentro dos limites.
\end{block}

\textbf{Método: GET}

\begin{block}{Request}
\begin{verbatim}
GET /caldeira/status HTTP/1.1
Host: api.siderurgia.com
Accept: application/json
\end{verbatim}
\end{block}

\begin{block}{Response}
\begin{verbatim}
HTTP/1.1 200 OK
Content-Type: application/json

{
  "temperatura": 850,
  "pressao": 120,
  "status": "OPERANDO_NORMAL"
}
\end{verbatim}
\end{block}

\end{frame}  
\begin{frame}
\centering
\vspace{1cm}

{\Huge \textbf{Estudo de Caso}} \\[0.5cm]
{\Large Teste de uso de sistema}

\vspace{1cm}

Sistema de Monitoramento de Caldeira Industrial \\
Siderurgia

\vfill
\small Exercício de Teste Manual de QA
\end{frame}


\begin{frame}[fragile]{GET – Consulta de Status da Caldeira}

\begin{block}{O que o QA fez}
O QA acessou o endpoint de status para verificar se a caldeira
está operando dentro dos limites.
\end{block}

\textbf{Método: GET}

\begin{block}{Request}
\begin{verbatim}
GET /caldeira/status HTTP/1.1
Host: api.siderurgia.com
Accept: application/json
\end{verbatim}
\end{block}

\begin{block}{Response}
\begin{verbatim}
HTTP/1.1 200 OK
Content-Type: application/json

{
  "temperatura": 850,
  "pressao": 120,
  "status": "OPERANDO_NORMAL"
}
\end{verbatim}
\end{block}

\end{frame}  
\begin{frame}[fragile]{POST – Tentativa com Dados Inválidos}

\begin{block}{O que o QA fez}
O QA tentou criar um registro de manutenção
com dados incompletos para testar validação.
\end{block}

\textbf{Método: POST}

\begin{block}{Request}
\begin{verbatim}
POST /caldeira/manutencao HTTP/1.1
Content-Type: application/json

{
  "tecnico": "",
  "prioridade": "ALTISSIMA"
}
\end{verbatim}
\end{block}

\begin{block}{Response}
\begin{verbatim}
HTTP/1.1 400 Bad Request

{
  "erro": "Prioridade inválida"
}
\end{verbatim}
\end{block}

\end{frame}
\begin{frame}[fragile]{POST – Criação de Manutenção Corretiva}

\begin{block}{O que o QA fez}
Após corrigir os dados, o QA criou corretamente
um novo registro de manutenção.
\end{block}

\textbf{Método: POST}

\begin{block}{Request}
\begin{verbatim}
POST /caldeira/manutencao HTTP/1.1
Content-Type: application/json

{
  "tecnico": "Ana Souza",
  "descricao": "Troca de sensor",
  "prioridade": "ALTA"
}
\end{verbatim}
\end{block}

\begin{block}{Response}
\begin{verbatim}
HTTP/1.1 201 Created
Location: /caldeira/manutencao/52
\end{verbatim}
\end{block}

\end{frame}
\begin{frame}[fragile]{GET – Consulta de Registro Inexistente}

\begin{block}{O que o QA fez}
O QA tentou consultar um registro inexistente
para validar tratamento de erro 404.
\end{block}

\textbf{Método: GET}

\begin{block}{Request}
\begin{verbatim}
GET /caldeira/manutencao/9999 HTTP/1.1
Host: api.siderurgia.com
\end{verbatim}
\end{block}

\begin{block}{Response}
\begin{verbatim}
HTTP/1.1 404 Not Found

{
  "erro": "Registro não encontrado"
}
\end{verbatim}
\end{block}

\end{frame}
\begin{frame}[fragile]{DELETE – Remoção de Registro Criado}

\begin{block}{O que o QA fez}
O QA removeu o registro recém-criado
para validar exclusão correta do recurso.
\end{block}

\textbf{Método: DELETE}

\begin{block}{Request}
\begin{verbatim}
DELETE /caldeira/manutencao/52 HTTP/1.1
Host: api.siderurgia.com
\end{verbatim}
\end{block}

\begin{block}{Response}
\begin{verbatim}
HTTP/1.1 204 No Content
\end{verbatim}
\end{block}

\end{frame}
\begin{frame}[fragile]{DELETE – Falha Interna no Servidor}

\begin{block}{O que o QA fez}
O QA tentou remover um registro durante
uma falha simulada no banco de dados.
\end{block}

\textbf{Método: DELETE}

\begin{block}{Request}
\begin{verbatim}
DELETE /caldeira/manutencao/60 HTTP/1.1
Host: api.siderurgia.com
\end{verbatim}
\end{block}

\begin{block}{Response}
\begin{verbatim}
HTTP/1.1 500 Internal Server Error

{
  "erro": "Falha ao acessar banco de dados"
}
\end{verbatim}
\end{block}

\end{frame}

\begin{frame}{Resposta – Classificação das Requisições}

\begin{block}{Resumo das Execuções do QA}
Durante o teste manual foram realizadas 6 requisições
utilizando os métodos GET, POST e DELETE.
\end{block} 

\begin{itemize}
    \item GET → 200 OK 
    \item POST → 400 Bad Request 
    \item POST → 201 Created 
    \item GET → 404 Not Found 
    \item DELETE → 204 No Content 
    \item DELETE → 500 Internal Server Error 
\end{itemize}

\begin{block}{Totais}
\begin{itemize}
    \item Total de requisições: 6
    \item Sucessos (2xx): 3
    \item Erros (4xx e 5xx): 3
\end{itemize}
\end{block}

\end{frame}
\begin{frame}{Resposta – Cálculo da Disponibilidade}

\begin{block}{Fórmula}
\[
Disponibilidade =
\frac{Requisições\ bem\ sucedidas}
{Total\ de\ requisições}
\]
\end{block}

\textbf{Abordagem 1 – Apenas 2xx como sucesso}

\[
Disponibilidade = \frac{3}{6} = 50\%
\]

\begin{block}{Interpretação}
Considerando somente códigos 2xx como sucesso,
a disponibilidade observada no teste foi de 50\%.
\end{block}

\end{frame}
\begin{frame}{Resposta – Interpretação Realista}

\begin{block}{Análise por Classe de Erro}
\begin{itemize}
    \item 4xx → erro causado pelo cliente
    \item 5xx → erro interno do servidor
\end{itemize}
\end{block}

Durante o teste houve apenas 1 erro 5xx.

\[
Disponibilidade = \frac{5}{6}
\]

\[
Disponibilidade \approx 83,33\%
\]

\begin{block}{Conclusão Técnica}
Se considerarmos apenas falhas do servidor (5xx),
a disponibilidade real do sistema foi de 83,33\%.

Erros 4xx indicam que o sistema respondeu corretamente
a uma requisição inválida.
\end{block}

\end{frame}
\end{document}
